\documentclass[conference]{IEEEtran}
\usepackage{cite}
\usepackage{amsmath,amssymb,amsfonts}
\usepackage{algorithmic}
\usepackage{graphicx}
\usepackage{textcomp}
\usepackage{xcolor}
\usepackage{listings}
\usepackage{booktabs}
\usepackage{hyperref}

\hypersetup{
    colorlinks=true,
    linkcolor=blue,
    filecolor=magenta,      
    urlcolor=cyan,
    pdftitle={Execution Identity for Agentic Systems: A Zero-Trust Runtime Governance Architecture},
    pdfpagemode=FullScreen,
}

\begin{document}

\title{Execution Identity for Agentic Systems: A Zero-Trust Runtime Governance Architecture}

\author{\IEEEauthorblockN{Jui-Cheng (Jimmy) Chen}
\IEEEauthorblockA{\textit{Top Celestial Company Ltd.} \\
Taipei, Taiwan R.O.C. \\
jimmychen@dr-os.io}
}

\maketitle

\begin{abstract}
As Multi-Agent Systems proliferate across enterprise and personal assistant applications, AI Agents are being widely deployed for retrieval and reasoning tasks in highly confidential environments. However, current mainstream frameworks universally rely on fragile ``prompt-based soft boundaries'' to control permissions, leading to a critical ``Agentic Trust Paradox.'' Because Large Language Models inherently suffer from Prompt Injection and other incurable vulnerabilities, once an agent is compromised, it can easily forge identities and trigger unauthorized access with catastrophic security consequences. To contain such threats and establish accountability mechanisms, this paper proposes the DROS (Deterministic Runtime Operating System) architecture. Unlike traditional application-level Workload Identity, this research introduces a finer-grained ``Execution Identity'' model. Through dynamically bound, By-Execution digital certificates issued by the Agent Identity Authority (AIA), combined with the GuardVM zero-trust gateway, DROS successfully intercepted all attacks in our prototype evaluation. The architecture achieves cryptographic Non-Repudiation traceability with minimal latency overhead, establishing an infrastructure standard for the emerging secure ``AgentWeb'' era.
\end{abstract}

\begin{IEEEkeywords}
Multi-Agent Systems, Zero-Trust, Execution Identity, By-Execution Certificate, Agent Security, Non-Repudiation, EU AI Act
\end{IEEEkeywords}

\section{Introduction}
In the process of enterprise adoption of generative AI and Retrieval-Augmented Generation (RAG), Multi-Agent collaboration has become the mainstream approach for handling complex cross-departmental operations. However, existing open-source frameworks are often built on an implicit assumption of ``complete internal trust.'' This architecture gives rise to a fatal ``Agentic Trust Paradox'': in enterprise systems that routinely chain dozens of Agents, how can a high-privilege decision-making agent verify that the instructions it receives genuinely originate from a downstream client with legitimate database access rights, rather than from a compromised agent executing unauthorized operations after malicious injection?

In commercial scenarios such as financial cross-departmental compliance, medical records retrieval, or enterprise source code management, such unauthorized access would lead to catastrophic data breaches and compliance disasters. To definitively close this vulnerability, DROS (Deterministic Runtime Operating System) aims to establish an Agentic Security Governance system that ensures every Agent invocation is backed by ``Determinism'' and ``cryptographic proof.''

\emph{``Unlike orchestration frameworks (e.g., LangChain, AutoGen) and communication protocols (e.g., MCP), DROS focuses on runtime identity, authorization, and non-repudiable governance. Therefore, DROS is designed as a complementary security substrate rather than a replacement for existing Agent ecosystems.''}

\section{Threat Model \& Security Assumptions}
In the theoretical construction of this study, we adopt an Adversarial Realism assumption:
\begin{enumerate}
\item \textbf{Untrusted Brain Assumption}: We assume that the LLM brain responsible for generating inference intents is completely untrusted and can be compromised at any time. The security boundary must not rely on the LLM brain remaining sane.
\item \textbf{Choke Point Enforcement Principle}: DROS abandons the illusion of container-level full syscall interception in the kernel space. The architecture guarantees that as long as a sensitive call is initiated through the registered \texttt{ToolExecutor} interface, it will be strictly constrained by cryptographic identity at the FFI boundary, with zero possibility of bypass.
\end{enumerate}

\section{Execution Non-Repudiation Model}
We abstract the execution security lifecycle of agents as a deterministic mathematical function evaluation process. Given an execution tuple:
\begin{equation}
\text{Decision} = f(\text{Policy}_{\text{Hash}}, \text{Agent}_{ID}, \text{Tool}_{\text{Name}}, \text{Context})
\end{equation}
To achieve absolute determinism and non-repudiability, this model strictly requires satisfying the following four abstraction bounds:
\begin{enumerate}
\item \textbf{Compile-time Resolution}: All conflict resolution and capability matrix expansion must be resolved and compiled into a binary capability bitmap (Capability Bitmap) at compile time.
\item \textbf{Constant-time Enforcement}: The time complexity of the security evaluation function $f$ at the runtime gate must be strictly constant-time $\mathcal{O}(1)$.
\item \textbf{No-Runtime-Intelligence}: The runtime interception sandbox must not parse strings, evaluate JSON text policies, or call any LLM for semantic decisions on the hot-path.
\item \textbf{Forensic Reproducibility}: Every execution trace must irreversibly bind the binary policy hash, the agent's identity signature, and the execution result.
\end{enumerate}

\section{System Architecture}
DROS operates as Runtime Security Substrate. AIA generates RSA-2048 key pair as Trust Anchor.
BEC binds execution\_id and allowed\_scopes. Single-use, eliminates credential theft and replay.
GuardVM plays a dual role: PDP + PEP.

\begin{figure}[htbp]
\centering
\includegraphics[width=3.5in]{DROS_Visual_Assets/dros_defense_layers_en.png}
\caption{DROS Runtime Security Layers and FFI/C-ABI Boundary Interception}
\label{fig:defense_layers}
\end{figure}

\subsection{Compile-Time: Deterministic Binary Capability Mapping}
The compiler performing lexicographical sorting mapping abstract \texttt{Agents} and \texttt{Tools} into indices. It generates a read-only flat file \texttt{policy.bin}.

\subsection{Runtime: FFI Interception and Lock-Free Bitwise Verification}
VajraClaw interposes at the C-ABI/FFI boundary. The authorization check is resolved to a low-level bitwise AND:
\begin{equation}
\text{Is\_Authorized} = (\text{Bitmap}[R_{\text{agent}}] \gg C_{\text{tool}}) \text{ \& } 0x01
\end{equation}

\begin{figure}[htbp]
\centering
\includegraphics[width=3.5in]{DROS_Visual_Assets/dros_vajra_concurrency_en.png}
\caption{Multi-Agent Capability Bitmap Concurrency Control}
\label{fig:concurrency}
\end{figure}

To avoid marshalling overheads, DROS uses zero-copy shared memory cyclic buffers.

\begin{figure}[htbp]
\centering
\includegraphics[width=3.5in]{DROS_Visual_Assets/dros_zerocopy_en.png}
\caption{Zero-Copy Shared Memory Pipeline between FFI Boundary and GuardVM}
\label{fig:zerocopy}
\end{figure}

\subsection{Agent PKI Network: Three-Tier Certificate Chain}
DROS implements a three-tier cryptographic certificate chain: Root CA (Trust Anchor), AIA (Issuing Authority), and BEC (Leaf Certificate).

\begin{figure}[htbp]
\centering
\includegraphics[width=3.5in]{DROS_Visual_Assets/dros_pki_chain_en.png}
\caption{DROS Agent Three-Tier Certificate Chain Key Management}
\label{fig:pki_chain}
\end{figure}

The AIA Dynamic Certificate Issuance and Revocation Lifecycle is illustrated below:

\begin{figure}[htbp]
\centering
\includegraphics[width=3.5in]{DROS_Visual_Assets/dros_lifecycle_flow_en.png}
\caption{AIA Dynamic Certificate Issuance and Revocation Lifecycle Flow}
\label{fig:lifecycle}
\end{figure}

\section{Empirical Evaluation}
To prove the physical correctness, we conducted stress testing using the DROS Home Lab verification harness on a PC, VM, and NAS.

\subsection{Experimental Setup \& Methodology}
\begin{itemize}
\item \textbf{Developer PC}: 8-core CPU, 16GB RAM client simulator.
\item \textbf{GuardVM Gateway}: Ubuntu Headless VM (4 vCPUs, 4GB RAM) running VajraClaw C-ABI native daemon.
\item \textbf{NAS Audit Vault}: Storage Vault connected via local network.
\item \textbf{Methodology}: Lookups run 100,000 times to eliminate OS scheduling jitter.
\end{itemize}

\subsection{Attribution \& Isolation Tests}
Under eBPF, the kernel only saw process \texttt{python3} executing syscalls, failing to isolate different Agent scopes. DROS intercepted FFI at C-ABI, checked the Capability Bitmap in \textbf{484 ns}, and killed the compromised thread immediately while letting legitimate administrator calls proceed normally.

\subsection{Revocation \& Concurrency Tests}
We simulated 50 concurrent agent threads with policy OTA swaps pushed every \textbf{10 ms}. DROS successfully processed \textbf{61,919 random tool calls} and completed \textbf{187 atomic policy updates} in memory with 0 race conditions.

\begin{table}[htbp]
\caption{DROS Core Latency Benchmarks}
\centering
\begin{tabular}{lll}
\toprule
Dimension & Measured Value & Significance \\
\midrule
Single Lookup Latency & 352.0 ns & Negligible in practical I/O \\
Peak Stress Latency & 484.8 ns & Sub-microsecond defensive stiffness \\
Time Complexity & Strictly $\mathcal{O}(1)$ & Flat latency curve w.r.t rule count \\
Runtime String Overhead & 0.000000 Alloc. & Zero JSON parser footprint \\
\bottomrule
\end{tabular}
\label{tab:perf}
\end{table}

\begin{figure}[htbp]
\centering
\includegraphics[width=3.3in]{DROS_Visual_Assets/dros_vs_opa_benchmark.png}
\caption{DROS vs OPA Latency Curve under Scale Limits}
\label{fig:vs_opa}
\end{figure}

\begin{figure}[htbp]
\centering
\includegraphics[width=3.3in]{DROS_Visual_Assets/dros_concurrency_latency.png}
\caption{DROS vs OPA Latency Curve under Thread Concurrency}
\label{fig:concurrency_latency}
\end{figure}

\begin{figure}[htbp]
\centering
\includegraphics[width=3.5in]{DROS_Visual_Assets/dros_rcu_en.png}
\caption{Lock-free Policy Hot-swap Concurrency Model based on RCU}
\label{fig:rcu_model}
\end{figure}

\subsection{Bypass \& Forgery Tests}
When integrated with OpenClaw framework, GuardVM blocked unauthorized direct FFI attempts. Expired/forged BECs failed validation at GuardVM, triggering immediate fail-closed shutdown.

\subsection{Non-Repudiation Audit Tests}
Logs securely shipped to NAS Audit Vault bind Agent ID and Policy SHA-256 Hash, providing legally binding forensic proof.

\section{Related Work}
\begin{enumerate}
\item \textbf{vs. SPIFFE/SPIRE}: SPIFFE stops at authentication and cannot enforce blocks at the syscall level.
\item \textbf{vs. OPA \& Rego}: OPA relies on dynamic JSON parsing with millisecond-level latency. DROS flattens rules into bitmaps, reducing lookup to 352 nanoseconds.
\item \textbf{vs. eBPF Policy}: eBPF sandboxes operate in the kernel and are context-blind to multi-agent contexts. DROS acts as the Identity Enforcement Layer.
\end{enumerate}

\section{Conclusion \& Future Work}
This study defined the Agent Runtime Attribution Gap and proved the feasibility of executing zero-trust microkernel interception. Future work will bind the Agent PKI with TEEs.

\section*{AI Collaboration Declaration}
During the writing of this paper and development of the prototype system, the authors utilized Large Language Models to assist with academic literature translation, code framework generation, and English formatting and proofreading. The authors assume full responsibility for all logical architecture, core algorithm design, and final content.

\begin{thebibliography}{10}

\bibitem{nist-zt}
S. Rose, O. Borchert, S. Mitchell, and S. Connelly, ``Zero Trust Architecture,'' \emph{NIST Special Publication 800-207}, Aug. 2020.

\bibitem{rsa}
R. L. Rivest, A. Shamir, and L. Adleman, ``A method for obtaining digital signatures and public-key cryptosystems,'' \emph{Communications of the ACM}, vol. 21, no. 2, pp. 120--126, 1978.

\bibitem{jwt}
M. Jones, J. Bradley, and N. Sakimura, ``JSON Web Token (JWT),'' \emph{RFC 7519}, May 2015.

\bibitem{aiact}
European Parliament and Council of the European Union, ``Regulation (EU) 2024/1689 of the European Parliament and of the Council of 13 June 2024 laying down harmonised rules on artificial intelligence (Artificial Intelligence Act),'' \emph{Official Journal of the European Union}, Jun. 2024.

\bibitem{mas}
M. Wooldridge, \emph{An Introduction to Multiagent Systems}, 2nd ed. Wiley, 2009.

\bibitem{spiffe}
E. Feldman et al., ``SPIFFE: Secure Production Identity Framework for Everyone,'' \emph{Cloud Native Computing Foundation (CNCF)}, 2024.

\bibitem{owasp-llm}
OWASP Foundation, ``OWASP Top 10 for Large Language Model Applications,'' \emph{OWASP Foundation}, 2025.

\end{thebibliography}

\end{document}
